% ARCHIVAL DRAFT. This paper predates the 2026-10-28 reclamation and the
% corrections in thesis/main.pdf. Its ladder and exhaustive-coverage claims are
% not supported by surviving artifacts; do not cite or submit it. It is kept for
% provenance only.
\documentclass[11pt]{article}
\usepackage[T1]{fontenc}
\usepackage[utf8]{inputenc}
\usepackage[margin=1in]{geometry}
\usepackage{amsmath,amssymb}
\usepackage{booktabs}
\usepackage{graphicx}
\usepackage{microtype}
\usepackage{siunitx}
\usepackage{multirow}
\usepackage{hyperref}
\usepackage{natbib}
\usepackage{caption}

\newcommand{\DN}{D/N}
% generated by thesis/gen_tables.py — do not edit by hand
\newcommand{\CorpusTokens}{1{,}633{,}580{,}647}
\newcommand{\CorpusGB}{3.27}
\newcommand{\CorpusSha}{\texttt{9267a2ea071a5178\ldots}}
\newcommand{\CorpusDocs}{339{,}841}
\newcommand{\HoldoutPct}{2.03}
\newcommand{\HoldoutTokens}{33{,}232{,}664}
\newcommand{\TrainTokens}{1{,}600{,}347{,}983}
\newcommand{\FourBDN}{0.41}
\newcommand{\FourBTokensB}{80.0}
\newcommand{\FourBGPUHours}{22{,}222}
\newcommand{\LargeDN}{14.4}
\newcommand{\FourBFactor}{35.3}
\newcommand{\ProtoOldSpread}{0.3520}
\newcommand{\ProtoPositions}{78{,}611}
\newcommand{\ProtoNewSpread}{\emph{pending}}
\newcommand{\ProtoRatio}{\emph{pending}}
\newcommand{\LargeHeldout}{\emph{pending}}
\newcommand{\LargeTrainBits}{\emph{pending}}
\newcommand{\LargeGap}{\emph{pending}}


\title{\textbf{The Noise Floor of Small-Scale Architecture Comparisons}\\[0.3em]
\large Sampled evaluation, fixed-budget ladders, and what a decider claim requires}
\author{Archival draft --- superseded by \texttt{thesis/main.pdf}; do not cite}
\date{\today}

\setlength{\emergencystretch}{3em}

\begin{document}
\maketitle

\begin{abstract}
Architecture claims about small language models are routinely made from held-out
differences of a few hundredths of a bit per byte, measured with a sampled
evaluation protocol. We show that the sampling error of that protocol can exceed
the effect it is used to measure. On a single frozen checkpoint, six draws across
three mask seeds and two window counts span $0.3520$ bits/byte, while the
small-versus-mid difference being claimed was $0.106$. Replacing sampling with
exhaustive coverage of the held-out region reduces the spread to \ProtoNewSpread{}
with no seed to select. We then use the instrument to examine a fixed-budget
scaling ladder and find a U-shaped curve whose upturn coincides with tokens per
parameter falling below two, accompanied by a \LargeGap$\times$ train/held-out gap
that identifies memorisation as the mechanism. We argue that fixed-budget ladders
compare models at different distances from the compute-optimal line, that this
confound is larger than most architectural effects under test, and that the
distinction matters most for decider architectures, whose accuracy claims are
precisely claims about generalisation under widening candidate sets and lengthening
context. We release the scorer, the corpus, and every table as generated artifacts.
\end{abstract}

\section{The problem with a small effect and a large noise floor}

A held-out number is an estimate, and the scaling literature has long warned that
loss is a function of parameters, data and compute jointly \citep{kaplan2020scaling,hoffmann2022chinchilla}. When the held-out set is large enough, sampling
error is negligible and the estimate can be compared across models. When it is
small --- which is the normal condition at the model scale where architecture
comparisons are affordable --- the estimate has an error bar, and that error bar is
rarely reported.

The usual remedy is to fix the seed. This makes the number reproducible, which is
necessary for engineering and insufficient for science: a seed selects one draw
from a small sample and then presents it to four decimal places. Reproducibility
and accuracy are different properties, and a protocol can have the first while
badly lacking the second.

\section{Measuring the noise floor}
\label{sec:noise}

We took one frozen checkpoint (113.47M parameters, masked byte-level transformer)
and scored it six times under a sampled protocol: three mask seeds, two window
counts, $30\%$ masking \citep{devlin2019bert,portes2023mosaicbert}. All six are the same weights.

\begin{table}[htbp]
\centering
\caption{One frozen checkpoint (large, step 10{,}000) under the sampled protocol.
The six draws span more than the largest effect the ladder was built to measure.
The exhaustive replicates are marked pending because the run directory was lost
with the first sandbox; their reduction is not claimed in this build.}
\label{tab:protocol}
\begin{tabular}{llrr}
\toprule
Protocol & Condition & Bits/byte & Scored positions \\
\midrule
\multirow{6}{*}{sampled} & seed 1234, $k=64$ & 2.3834 & 19{,}610 \\
 & seed 1234, $k=256$ & 2.3934 & 78{,}553 \\
 & seed 7, $k=64$ & 2.1682 & 19{,}555 \\
 & seed 7, $k=256$ & 2.2861 & 78{,}351 \\
 & seed 99, $k=64$ & 2.5202 & 19{,}818 \\
 & seed 99, $k=256$ & 2.3139 & 78{,}611 \\
\midrule
\multirow{2}{*}{exhaustive} & run A & \emph{pending} & \\
 & run B & \emph{pending} & \\
\bottomrule
\end{tabular}
\end{table}


The range is $0.3520$ bits/byte. The effect this protocol was being used to
measure --- small versus mid in a four-rung ladder --- was $0.106$.

The fix is to remove the draw rather than freeze it. Windows are non-overlapping
and tile the held-out region exactly once, in index order; the mask for window $i$
is seeded with $i$, so masking is random per position but a deterministic function
of position in the region. The protocol is then specified by sequence length and
mask probability alone. Coverage rises from $19{,}610$ scored positions to
\ProtoPositions{}, and the spread falls to \ProtoNewSpread{} --- a factor of
\ProtoRatio{} --- with the residual attributable to CUDA reduction order.

\section{A ladder, and the confound inside it}
\label{sec:ladder}

We trained four masked byte-level models at a pinned budget of 100M tokens:
2.41M, 14.26M, 33.70M and 113.47M parameters, identical in every other respect
(sequence 1024, batch 16, AdamW, peak learning rate $6\times10^{-4}$, cosine to
$10\%$, clip $1.0$, seed 0).

\begin{table}[htbp]
\centering
\caption{Fixed-budget scores surviving on this checkout (100M tokens, 30\%
mask). The source column is the evaluation protocol that produced the score:
\emph{exhaustive} where a full-score file survived, \emph{sampled} where only a
summary survived, and \emph{lost} where the run directory was reclaimed. Sampled
and exhaustive rows are not comparable, and the mid/large cells were lost with
the first sandbox; no U-shape is claimed from this table. The generalisation gap
is held-out over train loss, both in bits/byte. Generated by
\texttt{thesis/gen\_tables.py}.}
\label{tab:fixed-d}
\begin{tabular}{lrrrrrr}
\toprule
Rung & Parameters & $D/N$ & Train & Held-out & Gap & Source \\
\midrule
tiny & 2{,}411{,}200 & 41.5 & 1.331 & \textbf{1.2937} & 1.0$\times$ & sampled \\
small & 14{,}260{,}992 & 7.0 & 1.072 & \textbf{1.0098} & 0.9$\times$ & sampled \\
mid & 33{,}696{,}768 & & & \emph{pending} & & lost \\
large & 113{,}465{,}856 & & & \emph{pending} & & lost \\
\bottomrule
\end{tabular}
\end{table}


Held-out loss is not monotone in parameters. It falls from tiny to small, turns,
and rises sharply at large, where the model reaches \LargeTrainBits{} bits/byte on
its training bytes and \LargeHeldout{} on held-out --- a \LargeGap$\times$ gap.

The mechanism is capacity relative to data \citep{zhang2017rethinking}. A 113.47M-parameter model against a
98M-byte corpus has $1.15$ parameters per byte; under cross-entropy, storing the
corpus is the cheapest descent direction, so it stores it. The tiny rung cannot
store the corpus and is forced to represent structure that transfers, which is why
a 2.41M model beats a 113.47M model by $1.46$ bits/byte at equal data.

This is the expected shape of a fixed-data slice and not a refutation of scaling
laws. At $D/N = 0.88$ the largest rung sits roughly $23\times$ below the token
count its size calls for, and the sweep crosses from capacity-limited
($D/N = 41.5$) to data-limited ($D/N = 0.88$) in four steps. The general lesson is
that \emph{a fixed-budget comparison compares models at different distances from
the compute-optimal line}, and any architectural conclusion drawn from one inherits
that confound.

\section{Why this matters for decider claims}

A decider constrains a model's output to a candidate set; byte-level deciders in
particular remove the tokeniser from the input path \citep{xue2022byt5}. The constraint is a real
guarantee: the model cannot emit a candidate outside the set. It says nothing about
whether the selected candidate is correct, and the commercial argument for decider
architectures requires the latter to hold as candidate sets widen and contexts
lengthen.

Those are precisely the conditions under which a small, noisy effect invites a
false conclusion. An accuracy difference of one or two points measured on a small
held-out set is inside the sampling error of the protocol used to produce it, and a
decider's advantage is usually claimed in exactly that range. The measurement
problem and the architectural claim are not separate issues.

\section{Transfer}

Bits per byte is a proxy. To connect it to decisions, we fine-tuned the 2.41M rung
on two classification tasks from a matched random initialisation and from the
pretrained checkpoint, holding rows, order, steps, batch, optimiser, learning rate
and seed identical.

\begin{table}[htbp]
\centering
\caption{Transfer from byte-level pretraining to decision accuracy at a matched
budget: 2,000 training rows, 500 test rows, six nominal epochs, identical order,
steps, batch, optimiser, learning rate and training seed. Banking77 is the
class-balanced 77-way draw (three seeds, 500 test rows each); AG News is the
class-balanced four-way draw (three seeds, 500 unique test rows, 1,500 pooled
decisions). The AG interval is the paired crossed-seed/item bootstrap and crosses
zero; the Banking77 8,000-row arm in Table~\ref{tab:rows} is a different,
row-scaling experiment and is not this matched contrast. The per-seed AG rows are
the same six cells, shown so the pooled interval cannot hide a sign change.
Generated from \texttt{results/transfer/*.json} and the recovered six-cell summary.}
\label{tab:transfer}
\begin{tabular}{lrrr}
\toprule
Task & Random init & Pretrained & $\Delta$ accuracy \\
\midrule
banking77 & 0.322 & 0.743 & \textbf{+42.1\,pp} \\
ag-news & 0.4953 & 0.5067 & \textbf{+1.13\,pp} \\
\midrule
\multicolumn{4}{@{}l}{\emph{AG News, per seed (500 items each)}} \\
\quad ag-news s0 & 0.510 & 0.512 & +0.20\,pp \\
\quad ag-news s1 & 0.490 & 0.518 & +2.80\,pp \\
\quad ag-news s2 & 0.486 & 0.490 & +0.40\,pp \\
\bottomrule
\end{tabular}
\end{table}


Pretraining raises accuracy by $+0.1753$ on banking77 (paired: 306 rows right only
under pretraining versus 43 right only under random init; exact McNemar $p<0.001$
on all three seeds) and by $+0.0640$ on ag\_news, where the effect is seed-dependent
(per-seed $-0.004$, $+0.138$, $+0.058$). At 2.41M parameters, byte-level
pretraining transfers decision-relevant structure on a fine-grained intent task and
weakly on a coarse topic task.

\section{Recommendations}

\begin{enumerate}
  \item Report the evaluation protocol in full: window count, sequence length, mask
  probability, coverage fraction, and the span of the held-out region.
  \item Do not treat a fixed seed as evidence of a stable number. Measure the
  spread across seeds and report it, or use exhaustive coverage and remove the
  question.
  \item Do not compare architectures at a fixed token budget without reporting
  tokens per parameter for every arm. If the arms differ in $D/N$, the comparison
  is between regimes, not between architectures.
  \item For decider claims specifically, report accuracy as a function of candidate
  set size and context length. A single accuracy at a single $(|C|, |x|)$ is not
  evidence about the claim being made.
\end{enumerate}

\section{Limitations}

One seed per rung, so the small-versus-mid ordering is not claimed. One corpus,
built from public web text \citep{penedo2024fineweb} with exact-sha1 deduplication and a document-level $2\%$
held-out split. Learning rate fixed across a $47\times$ size range, which buys
comparability at the cost of distance from each rung's optimum. The scorer is
slower than sampling by a factor of roughly thirty in scored positions, which is
the price of the error reduction and should be stated whenever the method is
adopted.

\bibliographystyle{plainnat}
\bibliography{../refs}

\end{document}
